\section{Scope and the constructive arithmetic setting}\label{sec:foundations}

The Bell numbers provide a particularly tractable test of how much apparent
analysis can be expressed in intuitionistic primitive recursive arithmetic.
Their defining objects are finite, and the analytic formulas used below have
effective, explicitly bounded remainders.  Our aim is to give proofs of the
statements, not merely to observe that their numerical instances can be tested.
The target inventory is the two pages specified in the question
\cite{mathworld,wikipedia}; attribution of an identity to these pages does not
make either page an axiom of our argument.

\subsection{What ``in \ipra'' means here}

We use first-order intuitionistic logic with equality, symbols for primitive
recursive functions, their defining equations, the usual elementary successor
axioms, and induction for quantifier-free formulas.  An equivalent presentation
may use Boolean characteristic functions and induction for equations.  There is
no induction for arbitrary quantified formulas, no excluded middle for
arbitrary propositions, and no quantification over arbitrary functions or real
numbers.  Free numerical parameters are allowed in each induction instance.

The paper has two kinds of assertions.
\begin{enumerate}
\item \emph{Exact finite assertions}: equations, inequalities, divisibility,
finite counting statements and bounded algorithms, written directly using
primitive recursive operations.
\item \emph{Represented analytic assertions}: abbreviations for families of
finite rational inequalities with specified errors and primitive recursive
precision bounds.  For example, equality of two explicitly constructed real
names means that their rational approximants differ by at most the sum of the
declared errors, at every precision.  An asymptotic equivalence includes a
primitive recursive threshold for each rational accuracy request.
\end{enumerate}
This convention does not assert that the first-order theory contains a complete
ordered field of real numbers.  It also does not identify a theorem proved in
Lean's ambient type theory with a syntactic derivation in \ipra.  The distinction
is maintained in Section~\ref{sec:formalization}.

\subsection{Finite objects and bounded induction}

We encode a length-$l$ list of entries smaller than $b$ by its base-$b$ numeral,
with its length recorded separately.  Take $b\ge2$; exceptional smaller bases
are replaced by $2$.  Reading an entry uses division and remainder.  A table
with variable row lengths can instead use nested pairing; either convention is
primitive recursive.  A bounded sum, bounded product, bounded maximum and
bounded search are defined by primitive recursion.  Integers are signed
natural-number codes, and rationals are pairs $(z,d)$ with $d>0$.  Reduction to
lowest terms is optional.  Comparisons are decided by cross multiplication.
Polynomial and matrix computations use finite coefficient and entry lists.

A bounded formula is replaced by its Boolean characteristic function.  For
example, ``all $j\le m$ satisfy $P(j)$'' becomes the equation
\[
 \prod_{j=0}^{m}\chi_P(j)=1.
\]
Thus an induction on a bounded assertion is an instance of quantifier-free
induction after this explicit translation.  Nested induction on fixed finite
tuples is allowed: it is a finite sequence of induction arguments, not induction
along a new transfinite ordering.  Induction on a table of length $m$ is made
simultaneous by coding the whole table invariant in one characteristic function.

For a finite injection or bijection we supply algorithms for the map and its
inverse, together with bounds on their searches.  The counting lemma itself is
proved by induction while scanning the bounded domain.  Likewise, finite
pigeonhole is proved by keeping the Boolean table of already seen values.
These observations justify the finite counting arguments below without using
an infinitary set-existence principle.

\begin{remark}[A limitation of the compilation claim]
It is not true that every computable construction, every terminating Lean
program, or every classical proof of a numerical identity automatically belongs
to \ipra.  In this paper the loops have displayed primitive recursive bounds,
and the invariants used in the arithmetic arguments are bounded.  The prose
proof is an arithmetic proof blueprint; a complete generated object-language
derivation is a separate formalization task, not a deliverable claimed to have
already been checked.
\end{remark}

\subsection{Elementary analytic names}\label{subsec:names}

Only fixed, explicitly given families of names occur.  There is no claimed
primitive recursive evaluator for arbitrary primitive recursive programs.
An error parameter $s\in\N$ means absolute error at most $2^{-s}$.
When a named quantity is supplied as an argument, its data also include a
certified rational range and, for a logarithm or reciprocal, a strictly
positive rational lower bound. One first evaluates the argument to a
predetermined precision and then the rational approximation program. A
derivative bound $M$ on the certified range makes the propagated error at
most $M$ times the input error, plus the program's own error. This explains
composition of the names without exact sign tests on arbitrary real names.

For $|x|\le A$, with $A$ a positive integer, let
\[
 E_M(x)=\sum_{j=0}^{M}\frac{x^j}{j!}.
\]
If $M+2\ge2A$, the absolute tail is bounded by
\begin{equation}\label{eq:exp-tail}
 \sum_{j>M}\frac{A^j}{j!}\le
 \frac{2A^{M+1}}{(M+1)!}.
\end{equation}
The inequality means that every finite tail satisfies the bound; it follows
from a geometric majorant because the successive ratios are at most $1/2$.
One completely explicit precision rule is obtained by starting at $K=2A$:
$A^K/K!\le A^{2A}\le2^{2A^2}$, so
$M=2A+2A^2+s+3$ is more than sufficient.  This also proves convergence and the
addition and multiplication formulas by comparing finite products and their
tails.  Negative exponents may alternatively use the reciprocal of a positive
exponential, whose lower bound is $1$.

For $x>0$ rational, first find, by bounded integer comparisons, $k\in\mathbb Z$
and $y\in[1,2]$ such that $x=2^k y$.  The magnitude of $k$ is bounded by the sum
of the numerator and denominator of $x$.  With $z=(y-1)/(y+1)$, define
\begin{equation}\label{eq:log-series}
 \log y=2\sum_{j=0}^{J-1}\frac{z^{2j+1}}{2j+1}+R_J,
 \qquad
 0\le R_J\le\frac{2z^{2J+1}}{(2J+1)(1-z^2)}.
\end{equation}
Here $0\le z\le1/3$.  The identity is obtained by integrating the finite
geometric identity for $2/(1-z^2)$, retaining its remainder; equivalently it
can be proved by finite differentiation of the displayed polynomials.  Set
$\log x=k\log2+\log y$, using the same formula for $\log2$.  Taking
$J=s+|k|+4$ suffices for the present absolute-error convention.  The derivative
$1/x$, monotonicity, the product law and the inverse relation with $\exp$ follow
by finite polynomial approximation on an interval with an explicit positive
lower endpoint.  The error in replacing such an identity by a rational one is
bounded by the displayed tails and a bound for the derivative on that interval.

Square roots of nonnegative rational numbers are obtained by bisection using
the rational test $q^2\le x$.  Positive $d$th roots use $q^d\le x$ in exactly the
same way.  An initial interval $[0,1+\lceil x\rceil]$ and
$s+\lceil x\rceil+2$ bisections suffice.  Powers $x^\alpha$ at positive named
$x$ mean $\exp(\alpha\log x)$ unless a root construction is simpler.

The sine and cosine names are their factorial series.  The same absolute-tail
estimate as \eqref{eq:exp-tail} applies on a bounded interval.  One may define
$\pi$ by the convergent arctangent construction
$\pi=16\arctan(1/5)-4\arctan(1/239)$, where
\[
 \arctan t=\sum_{j=0}^{J-1}\frac{(-1)^j t^{2j+1}}{2j+1}
             +\rho_J,\qquad |\rho_J|\le\frac{|t|^{2J+1}}{2J+1}.
\]
The addition formula for tangent verifies this normalization; it agrees with
$4\arctan1$, whose alternating partial sums also give a slower primitive
recursive name.  Polynomial approximation proves the addition formulas for
sine and cosine, their derivatives, and $\exp(it)=\cos t+i\sin t$.
In particular, orthogonality of integer Fourier modes follows by integrating
their explicit sine and cosine antiderivatives.

\subsection{Integration and Taylor estimates as finite statements}

For the functions used here, an integral over a rational compact interval is a
name given by equal-mesh Riemann sums.  If $|f'|\le M$ on an interval of length
$L$, the error of a left-endpoint sum with $q$ subintervals is at most
$ML^2/q$.  Indeed, on each subinterval the oscillation is at most its length
times $M$; summing the $q$ resulting rectangles gives the bound.  The usual
derivative bounds used below are obtained directly from factorial series.
Endpoints such as $\pi$ are approximated with a certified error, adding at most
$\sup|f|$ times the endpoint error.  Thus all integrations in the proofs have
primitive recursive mesh bounds.

Taylor's estimate
\begin{equation}\label{eq:taylor-estimate}
 \left|f(x+h)-\sum_{j=0}^{d}\frac{f^{(j)}(x)}{j!}h^j\right|
 \le\frac{M|h|^{d+1}}{(d+1)!}
\end{equation}
holds when $|f^{(d+1)}|\le M$ on the segment.  A constructive proof integrates
the derivative identity repeatedly and bounds the resulting simplex integral;
the simplex has volume $|h|^{d+1}/(d+1)!$, verified by induction.  Equivalently,
one first proves the formula for polynomials and then uses uniform factorial
tails.  No intermediate-value theorem or selection of an unknown Taylor point
is used.  The elementary estimates $1-u\le e^{-u}$,
$|e^z-1|\le e^{|z|}-1$, and $1-\cos t\ge t^2/3$ for $|t|\le1$ follow from
these same finite series estimates.

Unbounded Gaussian integrals are notation for truncated integrals with the
tail certificate
\[
 \int_T^\infty e^{-t^2/2}\,dt\le T^{-1}e^{-T^2/2}\quad(T>0).
\]
Integration by parts gives
\[
 \frac1{\sqrt{2\pi}}\int_{-\infty}^{\infty}t^{2j}e^{-t^2/2}\,dt
 =(2j-1)!!,
 \qquad \int_{-\infty}^{\infty}t^{2j+1}e^{-t^2/2}\,dt=0.
\]
For completeness, the normalizing integral is computed by squaring a truncated
Gaussian integral, sandwiching the square between inscribed and circumscribed
disks, and using finite radial annuli. Integration of
$r e^{-r^2/2}$ gives the limit $2\pi$ for the square, with errors bounded by
Gaussian tails. More explicitly, an inscribed regular $q$-gon has area
$(q/2)R^2\sin(2\pi/q)$; for $q\ge8$, its difference from $\pi R^2$ is at
most $4\pi^3R^2/(3q^2)$ by the cubic sine remainder. The circumscribed
polygon has area $qR^2\tan(\pi/q)$, with excess at most
$2\pi^3R^2/(3q^2)$: use
$\sin x-x\cos x=\int_0^x t\sin t\,dt\le x^3/3$ and
$\cos(\pi/q)\ge1/2$. These two
polygonal computations first establish the disk-area name. On radial steps
$R/M$, the function $e^{-r^2/2}$ oscillates by at most $R^2/M$; the total
annular integration error is at most $4R^4/M$. Taking $q,M$ to be explicit
powers of $2$, with the requested error and $R$ as parameters, therefore
justifies the disk integral calculation by finite polygonal and rectangular
sums. This is an effective calculation, not an appeal to general measure
theory.

Every passage to a limit below can consequently be read as follows: fix a
rational requested error, choose the indicated finite truncations and meshes,
prove the finite rational inequality, and add its explicit error allowances.
The exact integer conclusions are unaffected by this notation.
